\chapter{Perforated Eardrum}

\section*{Definition}
A perforated eardrum is a full-thickness defect in the tympanic membrane separating the external auditory canal from the middle ear; perforations may be acute (traumatic or infectious) or chronic, and are classified by size, location, and duration.

\section*{What Happens in Your Body}
Disruption of the tympanic membrane compromises its barrier and sound-transmission function, causing conductive hearing loss proportional to the size of the perforation and exposing the middle ear to pathogens, water, and pressure changes.

\section*{Causes}
\begin{itemize}
  \item \textbf{Acute otitis media:} Purulent middle ear effusion causing pressure necrosis and spontaneous rupture
  \item \textbf{Barotrauma:} Rapid pressure change (explosion, slap, diving, airplane descent)
  \item \textbf{Trauma:} Direct injury from cotton swab, hairpin, or foreign body
  \item \textbf{Thermal:} Welding spark or slag entering the ear canal
  \item \textbf{Chronic infection:} Chronic suppurative otitis media with persistent or recurrent perforation
  \item \textbf{Iatrogenic:} Post-tympanostomy tube extrusion or myringotomy site
\end{itemize}

\section*{When to See a Doctor}
See your doctor if:
\begin{itemize}
  \item Sudden ear pain with immediate relief followed by otorrhea (suspect acute perforation)
  \item Bloody or purulent discharge from the ear
  \item Hearing loss after trauma or infection
  \item Tinnitus or vertigo after head trauma (suspect ossicular disruption or inner ear injury)
  \item Perforation persisting beyond 3 months (requires ENT evaluation)
\end{itemize}

\section*{Related Symptoms}
\begin{itemize}
  \item Conductive hearing loss (typically 20--30 dB)
  \item Otorrhea (bloody or purulent)
  \item Tinnitus
  \item Aural fullness
  \item Vertigo (if inner ear involved or cholesteatoma present)
\end{itemize}
\textbf{Important:} Chronic perforation with persistent purulent otorrhea defines chronic suppurative otitis media (CSOM) and requires aural toilet and antibiotic drops; cholesteatoma must be excluded.

\section*{Self-Care / Home Management}
\begin{itemize}
  \item Keep the ear absolutely dry (no swimming, showering with a cotton ball coated in petroleum jelly)
  \item Avoid blowing the nose forcefully (can force infection into the middle ear)
  \item No ear drops of any kind unless prescribed by a physician
  \item Over-the-counter pain relievers for associated discomfort
  \item Avoid air travel and scuba diving until perforation heals
  \item Avoid inserting any object into the ear canal
\end{itemize}

\section*{How Your Doctor Will Evaluate You}
\begin{itemize}
  \item Otoscopy: Visualization of the perforation margin; assess size, location (central vs. marginal), and presence of cholesteatoma
  \item Pneumatic otoscopy: TM immobility
  \item Tuning fork tests (Rinne: Air-bone gap; Weber: Lateralizes to affected ear)
  \item Pure-tone audiometry to quantify conductive loss and assess sensorineural function
  \item Otomicroscopy for detailed assessment of the perforation and middle ear mucosa
  \item CT temporal bone if cholesteatoma, ossicular erosion, or labyrinthine fistula suspected
\end{itemize}

\section*{Treatment Options}
\begin{itemize}
  \item \textbf{Acute small perforation:} Observation for 4--6 weeks (85--90\% heal spontaneously)
  \item \textbf{Infected perforation:} Topical antibiotic drops (fluoroquinolone) after aural toilet
  \item \textbf{Patch technique:} Paper patch or fat myringoplasty for small persistent perforations (offices procedure)
  \item \textbf{Tympanoplasty:} Surgical repair using temporalis fascia, perichondrium, or cartilage graft for perforations > 3 months
  \item \textbf{Underlay technique:} Graft placed medial to the TM remnant (standard approach)
  \item \textbf{Overlay technique:} Graft placed lateral to the fibrous layer (for large or anterior perforations)
  \item \textbf{Ossicular chain reconstruction:} Concurrent with tympanoplasty if ossicular erosion present
\end{itemize}

\section*{References}
\begin{thebibliography}{99}
\bibitem{ref1} Ott MC, Lundy LB. "Tympanic membrane perforation in adults." \textit{American Family Physician}, 2011;83(5):557--562.
\bibitem{ref2} Sheehy JL, Glasscock ME. "Tympanic membrane perforation: diagnosis and management." \textit{Laryngoscope}, 2002;112(3):535--543.
\bibitem{ref3} Schraff SA, Markham JL, Seibert JP, et al. "Traumatic tympanic membrane perforation: a review of 80 cases." \textit{Journal of Emergency Medicine}, 2020;58(1):78--84.
\bibitem{ref4} Semaan MT, Megerian CA. "The pathophysiology of cholesteatoma." \textit{Otolaryngologic Clinics of North America}, 2006;39(6):1043--1059.
\bibitem{ref5} Stangerup SE, Tixier P, Tos M, et al. "The natural history of spontaneous tympanic membrane perforation." \textit{Otology \& Neurotology}, 2003;24(4):561--564.
\bibitem{ref6} Zwierzchowska A, Gawlik J, Wiatr M, et al. "Tympanoplasty: current concepts and long-term outcomes." \textit{Otolaryngologia Polska}, 2019;73(6):1--7.

\end{thebibliography}

\clearpage
